\section{连接后续研究线索的三个跟算专题}
本节把新增文献所需的模型和推导仍集中在第一章。后面的研究脉络只讨论这些模型带来的发现、检验和限制。三个例子分别说明：一个闭合为何能改变电流层；一个几何比例为何能改变整体率；一个轨道不稳定指标为何不能直接替代温度和熵。

\subsection{俘获电子状态方程怎样支撑延长电流层}\label{sec:synthesis-closure-model}
\paragraph{第一步：分清演化方程与闭合。}
\citet{F611}的离子由粒子推进，电子无质量、准中性，其电场来自
\begin{equation}
\E=-\vct u_i\times\B+\frac{\J\times\B}{en}
-\frac{\grad\cdot\vct P_e}{en}
+\eta\J-\eta_H\nabla^2\J .
\end{equation}
电子压力取旋转对称形式
$\vct P_e=p_{e\perp}\vct I+(p_{e\parallel}-p_{e\perp})\vct b\vct b$，
$\vct b=\B/B$。这里没有演化六个独立压力分量，而是用$n,B$及其上游参考值求两个压强。因此它既不是等温Hall模型，也不是完整十矩。

\paragraph{第二步：写出闭合，并检查它到底近似了什么。}
定义$\widetilde n=n/n_\infty$、$\widetilde B=B/B_\infty$、
$\widetilde p=p/p_\infty$，以及
$a=\widetilde n^3/\widetilde B^2$。本文把原文的$\alpha$改记$a$，以免与后面的轨道参数混用。所实现的近似状态方程为
\begin{align}
\widetilde p_{e\parallel}
 &=\widetilde n\frac{2}{2+a}
   +\frac{\pi\widetilde n^3}{6\widetilde B^2}
      \frac{2a}{2a+1},\label{eq:synthesis-eos-par}\\
\widetilde p_{e\perp}
 &=\frac{\widetilde n}{1+a}
   +\widetilde n\widetilde B\frac{a}{1+a}.
\end{align}
弱俘获侧的第一项给近等温响应；强俘获侧趋向
$p_{e\perp}\propto nB$与$p_{e\parallel}\propto n^3/B^2$。
后者的幂次可以从本章CGL推导再次理解：守恒磁矩给
$T_\perp\propto B$；一段物质磁管同时守恒粒子数$nA\ell$和磁通$BA$，
故$\ell\propto B/n$；平行绝热作用量给$v_\parallel\ell$近恒定，于是
$T_\parallel\propto(n/B)^2$，再乘$n$便得$p_\parallel\propto n^3/B^2$。
需要慢变、电子充分磁化和适用的俘获排序；这里的简单磁管推理解释极限幂次，\textbf{不等于从Vlasov完整导出了插值函数}。

式\eqref{eq:synthesis-eos-par}中的权重和$\pi/6$系数取自该文采用的闭合，不是由两端极限唯一确定。本次尚未核读其前序完整俘获分布推导，不能补写成已核实的精确结果。一个可以自行核对的小细节是：在$\widetilde n=\widetilde B=1$处，插值给
$\widetilde p_\parallel=2/3+\pi/9\simeq1.0157$，而
$\widetilde p_\perp=1$；它是近似插值，并非处处精确恢复归一条件的恒等式。

\paragraph{第三步：由压力张量求曲率力。}
令$\Delta p=p_\parallel-p_\perp$、$\vct\kappa=(\vct b\cdot\grad)\vct b$。
按乘积法则，
\begin{equation}
\grad\cdot\vct P
=\grad p_\perp+\vct b(\vct b\cdot\grad\Delta p)
 +\Delta p\,\vct\kappa+\Delta p\,\vct b(\grad\cdot\vct b).
\end{equation}
向垂直于$\vct b$的平面投影，两个沿$\vct b$的项消失，得到
\begin{equation}
(\grad\cdot\vct P)_\perp=\grad_\perp p_\perp+\Delta p\,\vct\kappa .
\end{equation}
无质量电子动量平衡为
$\vct J_e\times\B=en\E+\grad\cdot\vct P_e$。
左乘$\B\times$并使用
$\B\times(\vct J_e\times\B)=B^2\vct J_{e\perp}$，
可把压力驱动部分写成
\begin{equation}
\vct J_{e\perp}^{(P)}
=\frac{\B\times\grad\cdot\vct P_e}{B^2},\qquad
\vct J_{e\perp}^{(\mathrm{anis})}
=\frac{\Delta p_e}{B}\vct b\times\vct\kappa .
\end{equation}
这就是各向异性通过弯曲场线支撑额外电流的途径；电流再经Ampère定律改变面外磁场，后者又进入状态方程。反馈不要求总重联率同时大幅变化。

\paragraph{第四步：解释火管参数，但不把电子参数当整个等离子体的稳定性判据。}
在同一曲率方向比较压力与磁张力，其系数是
\begin{equation}
\frac{B^2}{\mu_0}-\Delta p_e
=\frac{B^2}{\mu_0}(1-F_e),\qquad
F_e=\frac{\mu_0\Delta p_e}{B^2}.
\end{equation}
$F_e$接近一表示电子各向异性与磁张力同量级。完整物种合力还含离子压力、惯性和几何项；不能把$F_e=1$当成任意动理学模的精确阈值。F611实际长层中常为0.5--0.7，并非处处严格等于一。

\paragraph{第五步：上游标签也是模型的一部分。}
Harris初值$B_x=B_0\tanh(z/\lambda)$可取
$\Psi=-B_0\lambda\ln\cosh(z/\lambda)$，从而
\begin{equation}
\exp\!\left(\frac{2\Psi}{B_0\lambda}\right)
=\operatorname{sech}^2(z/\lambda),\quad
n_\infty(\Psi)=n_b+n_0\exp\!\left(\frac{2\Psi}{B_0\lambda}\right).
\end{equation}
这里的第二式只是从初值重建参考密度的起点。作者发现直接使用会在中心产生非物理各向异性，又给\emph{参考密度}加上中心局域修正，使$a$变小；它没有改变实际粒子密度。这个处理依赖二维磁通标签和初值信息。在复杂三维重联中如何定义、更新这些上游量，正是闭合尚未解决的一部分。

\subsection{有限X线模型：定义恒等式、物理假设与拟合分别在哪里}\label{sec:synthesis-length-model}
\paragraph{第一步：从总磁通开始。}
在\citet{F792}的特殊对称几何中，以$z=0$中面重联法向场积分定义$\Psi$，令薄片沿电流方向长$L_y$，活跃部分长$L_a$。两种无量纲率定义为
\begin{equation}
R=\frac{\dot\Psi}{B_0\VA L_y},\qquad
R_*=\frac{\dot\Psi}{B_0\VA L_a}.
\end{equation}
只要$L_a>0$，立刻有
\begin{equation}
R=\frac{L_a}{L_y}R_* .
\label{eq:synthesis-active}
\end{equation}
到此尚无新动力学，这是分母不同导致的恒等关系。物理内容是模拟显示活跃区的$R_*$趋于相近量级，以及为什么$L_a<L_y$。

\paragraph{第二步：用输运时间估计受抑长度。}
电子把新生法向磁通沿X线搬运。在流出尚未建立、片尚未变薄的时间$\tau_{\rm thin}$内，
\begin{equation}
L_s\sim |u_{ey}|\tau_{\rm thin},\qquad L_a\simeq L_y-L_s.
\end{equation}
若$|u_{ey}|\sim\VA$、$\tau_{\rm thin}\sim10\Omega_i^{-1}$，由
$\VA/\Omega_i=d_i$得到$L_s\sim10d_i$。
速度非均匀时，应比较
$\tau_{\rm tr}=\int\dd y/|u_{ey}(y)|$和$\tau_{\rm thin}$，
而不是强行用单个速度乘长度。此解释依赖电子磁通输运和进一步变薄的反馈；它不是Hall项在X点直接打破电子冻结的推导。\citep{F659}

\paragraph{第三步：求错位的电流和磁场怎样削弱出流。}
简化离子动量预算后，原文以
$V_{x,\max}^2\propto(J_yB_z)_{\max}$近似出流驱动，并把压力梯度影响吸收进系数。这不是严格一般恒等式。再假设
$J_y=J_m f(y)$、$B_z=B_m f(y+L_s)$，
其中$f$为平台两侧线性过渡的梯形，各侧过渡宽$L_t$。
长平台仍有重叠时，$f f_{\rm shift}$最大值为1。
若只有两段斜坡重叠，取两斜坡各自的剩余距离$s$和$L_a-s$，则
\begin{equation}
f f_{\rm shift}=\frac{s(L_a-s)}{L_t^2}
=\frac{L_a^2/4-(s-L_a/2)^2}{L_t^2}.
\end{equation}
最大值在$s=L_a/2$，所以
\begin{equation}
V_{x,\max}\simeq V_m
\begin{cases}
1,& L_a\geq2L_t,\\
L_a/(2L_t),&0<L_a<2L_t.
\end{cases}
\end{equation}
图5实际取$V_m=0.75\VA$、$L_t=4d_i$、$R_*=0.15$，$L_s$也由本组模拟拟合。因此成功的是受参数约束的机制模型，而非从任意初值无参数预测出流。

\note{跟算例：仅为练习，设固定$L_s=10d_i$、$R_*=0.15$。$L_y=30d_i$时式\eqref{eq:synthesis-active}给$R=0.10$；$L_y=15d_i$时为0.05。相同局地处理能力，整条X线的平均率仍不同。$L_y\leq10d_i$时不能把公式继续代入得到负率，应停止该活跃区近似。原文最短例仍有种子驱动瞬态，不能用零稳态近似抹掉。}

\subsection{轨道发散、内能与熵：怎样核对Hall加热解释}\label{sec:synthesis-heating-model}
\paragraph{第一步：先守住总能量恒等式。}
在单电荷准中性体系，
$\E_H=\J\times\B/(en)$和$\J=en(\vct u_i-\vct u_e)$，所以
\begin{align}
\J\cdot\E_H&=0,\\
en\vct u_i\cdot\E_H-en\vct u_e\cdot\E_H
 &= (\vct u_i-\vct u_e)\cdot(\J\times\B)=0.
\end{align}
两物种分别的电功可以非零，和却为零。Hall作用可组织电场、磁张力与物种能量转移；最终内能还需在压力应变、热流和碰撞预算中追踪。这与MRX测得离子受到很强面内电场加速并不矛盾。\citep{Yoo2013,L272}

\paragraph{第二步：在可控的局部例子里推轨道分离。}
取静态均匀$\B=B_0\vct e_z$，电场$E_x(x)$；离子电荷$q>0$、质量$m$。
两条相邻轨道的差$\delta x,\delta y$满足
\begin{equation}
\delta\ddot x=\frac qm E_x'(x)\delta x+\Omega_c\delta\dot y,\qquad
\delta\ddot y=-\Omega_c\delta\dot x,\quad \Omega_c=\frac{qB_0}{m}.
\end{equation}
在\emph{局地梯度冻结为常数}的近似下，对第一式求导并代入第二式：
\begin{equation}
\delta\dddot x=\Omega_c^2(a_x-1)\delta\dot x,\qquad
a_x=\frac{mE_x'}{qB_0^2}.
\end{equation}
若$a_x>1$，局部指数$\lambda^2=\Omega_c^2(a_x-1)>0$。
这里$E_x'$带符号。若沿轨道变化明显，还应有
$(q/m)(\dd E_x'/\dd t)\delta x$，不能直接删去。
因此文献常用的
$m|\grad_\perp\cdot\E_\perp|/(qB^2)$是提示失磁化区域的指标，不是任意时空电磁场的充分混沌判据；在$B$接近零处更不能单独依赖这一比值。

\paragraph{第三步：复算倒抛物势，消除系数歧义。}
令$\phi=-E_0(x^2+y^2)/(2\ell)$、$z_c=x+\mathrm i y$。
直接代入Lorentz方程得
\begin{equation}
\ddot z_c=\frac{qE_0}{m\ell}z_c-\mathrm i\Omega_c\dot z_c.
\end{equation}
若定义$a_p=4mE_0/(q\ell B_0^2)$，则第一项系数应为
$\Omega_c^2a_p/4$。代入$z_c\propto e^{\mathrm i\omega t}$可得
\begin{equation}
\omega^2+\Omega_c\omega+\frac{\Omega_c^2a_p}{4}=0,\qquad
\frac{\omega}{\Omega_c}=\frac{-1\pm\sqrt{1-a_p}}{2}.
\end{equation}
这与\citet{L272}式20相符，但其式18、19按所印$a_p$定义缺少$1/4$。
本报告修正教学代数，不据此推断试粒子代码也错了。
$a_p>1$给线性发散螺旋；这个可显式求解的线性系统说明\textbf{指数分离本身不足以证明不可逆热化}。

\paragraph{第四步：计算真正的温度需要减去整体流。}
用分布函数定义
\begin{equation}
K=\int \frac{mv^2}{2}f\,\dd^3v,\quad
\frac32 nT=\int\frac m2|\vct v-\vct u|^2f\,\dd^3v
=K-\frac12mn u^2 .
\end{equation}
温度$T$在本报告按能量计。它与
$\sqrt{\langle\epsilon^2\rangle-\langle\epsilon\rangle^2}$不是同一个量，
后者还涉及速度四阶矩；L272试粒子图中的TE采用后者。
两者都可描述展宽，但不能不加说明地互换。
无碰撞光滑Vlasov输运仍保持细粒度熵；温度上升、粗粒化熵增和经弱碰撞产生的熵增，还需分别检验。这样才能把轨道候选机制与本章相混合/Hermite推导、以及第5章的光谱测温接起来。
